% Round 15. Conditional on the arbitrary-base P1 output theorem and its
% specified product base-change property. Not a claim about rank > 2 bundles.
\section{A global diagonal kernel for rank-two projective bundles}
\label{r15b:rank-two-section}

The construction in this section uses smooth metrics on two line
bundles, not a smooth norm on the rank-two vector bundle itself.
This distinction is substantive in non-archimedean geometry.
The local analytic output input is recorded explicitly before it
is used.

\begin{definition}[The local projective-line output input]
\label{r15b:HP}
Write $(\mathrm{HP})_K$ for the following assertions on every
ordinary open $U$ in a smooth algebraic $K$-variety.
\begin{enumerate}
\item The prescribed current pushforward for
$\operatorname{pr}:U\times\mathbf P^1\to U$ of every ordinary
PWS form of bidegree $(p,q)$ has an ordinary PWS representative
of degree $(p-1,q-1)$, with all ordinary smooth compact-test
pairings. Degrees with $p=0$ or $q=0$ give zero. The assertion
is local on $U$ and applies to support proper over paired
compact subsets.
\item For the fixed reference diagonal kernel of Round 14,
the three outputs on every PWS $(1,1)$ input on
$U\times\mathbf P^1$ have ordinary PWS representatives of
degrees $(0,0),(1,0),(0,1)$ with all prescribed smooth mixed-test
pairings.
\item The projection in the first item has the coordinate
base-change identity for
$U\times\mathbf P^1_y\to U$, and hence for a morphism locally
of this form. In particular an input pulled back from
$U\times\mathbf P^1_z$ has the pullback of its $z$-integral
as integral on $U\times\mathbf P^1_z\times\mathbf P^1_y$.
This is an identity of ordinary PWS tensors with all mixed-test
pairings, not merely an equality on exterior-product tests.
\end{enumerate}
\end{definition}

The new arbitrary-base projection and kernel construction is the
intended source of $(\mathrm{HP})_K$. The base-change clause is
part of the hypothesis until its separate tensor/finite-descent
proof is accepted. The results below are conditional on precisely
this input; they do not infer that clause from fiberwise scalar
integration alone.

\begin{lemma}[Smooth line metrics and their actual classes]
\label{r15b:smooth-line-metrics}
Let $M/K$ be smooth integral projective and let $L$ be an
algebraic line bundle on $M$. It admits an ordinary smooth
metric. Its curvature is a separately closed ordinary smooth
$(1,1)$ form. Relative to any model metric on the same line,
the difference of the two first Chern forms is $D$ of a global
ordinary PWS scalar. Thus their actual Bott--Chern and Aeppli
classes agree.
\end{lemma}

\begin{proof}
Choose an ordinary trivializing cover with local frames $e_j$
and an ordinary smooth partition of unity $\rho_j$ subordinate
to it. Define, for a local nonvanishing section $s$,
\[
                    \log\|s\|=\sum_j\rho_j\log|s/e_j|.
\tag{PB15.1}\label{r15b:smooth-metric-formula}
\]
Each summand is extended by zero off the support of $\rho_j$.
The support condition makes it smooth across that extension,
and the expression has the metric transformation law because
$\sum_j\rho_j=1$. Local weights are smooth; their $D$ values
glue because transition functions are analytic units and their
logarithms have zero $D$. Both actual differentials kill the
resulting curvature.

The ordinary partition and metric facts are also CLD,
Propositions 15.2.8 and 17.1.7--17.1.8 in the 2025 version.
The space here is good, Hausdorff and compact, so their
paracompactness and local separation hypotheses hold.
A model metric has local PL weights. Subtracting the two
weights gives a global PWS metric-ratio scalar; applying $D$
gives the curvature difference. This proves the class statement
in Mihatsch's actual theory. No positivity of either metric
is required.
\end{proof}

Let now $S/K$ be smooth integral projective, let $N$ be a
rank-two vector bundle, and use the convention of lines in
\[
                \pi:E=\mathbf P_S(N)\longrightarrow S.
\]
Choose smooth metrics on the line bundles
$\mathcal O_E(1)$ and $\det N$, with curvatures $h$ and $\kappa$,
respectively. Put $T=E\times_S E$, with projections $p_z,p_y$
and structural map $\pi_T:T\to S$.

\begin{lemma}[The global relative diagonal Green scalar]
\label{r15b:rank-two-kernel}
On $T$, the diagonal is the zero divisor of a canonical section
$s_\Delta$ of
\[
 \mathcal L_\Delta
   =p_z^*\mathcal O_E(1)\otimes p_y^*\mathcal O_E(1)
                         \otimes\pi_T^*\det N.
\tag{PB15.2}\label{r15b:diagonal-line}
\]
For the tensor product of the chosen metrics, the scalar
\[
                         g_E=-\log\|s_\Delta\|
\]
is ordinary PWS on $T^{\mathrm{an}}\setminus\Delta^{\mathrm{an}}$
and satisfies the scalar Green identity
\[
             Dg_E=h_z+h_y+\pi_T^*\kappa-\delta_\Delta.
\tag{PB15.3}\label{r15b:rank-two-Green}
\]
The determinant term has a plus sign in the convention of lines.
Locally on a trivializing Zariski open $U\subset S$, it has
the form
\[
          g_E=g_{\mathrm{ref}}+a(x,z)+a(x,y)+b(x),
\tag{PB15.4}\label{r15b:rank-two-local-kernel}
\]
where $a$ is an ordinary PWS scalar on $U\times\mathbf P^1$
and $b$ is an ordinary smooth scalar on $U$.
\end{lemma}

\begin{proof}
The tautological lines on the two factors are subbundles
$L_z,L_y\subset\pi_T^*N$. Their wedge defines
\[
                 L_z\otimes L_y\longrightarrow\pi_T^*\det N.
\]
It is a section of the line in \eqref{r15b:diagonal-line}.
In a local frame of $N$ it is the determinant of the two
homogeneous vectors. Its zero scheme is exactly the relative
diagonal, with multiplicity one. This proves both the divisor
and the sign of the determinant factor.

The scalar actual Poincare--Lelong formula for a metrized line
bundle gives \eqref{r15b:rank-two-Green}; its smooth metric
weights are actual PWS weights. On a trivializing open of $N$,
compare the chosen metric on $\mathcal O_E(1)$ with the fixed
reference metric on the projective line. Their ratio is the
PWS scalar $a$. The determinant frame has smooth negative log
norm $b$. The determinant section and the tensor metric then
give \eqref{r15b:rank-two-local-kernel} literally. No metric on
the rank-two bundle $N$ is used. In particular this does not
contradict the warning in CLD Remark 17.1.9 about smooth
ultrametric norms on vector bundles of rank greater than one.
\end{proof}

\begin{proposition}[Actual outputs for the global rank-two kernel]
\label{r15b:rank-two-outputs}
Assume $(\mathrm{HP})_K$. For every
$\alpha\in\operatorname{PS}^{1,1}(E)$, the currents
\[
 G_E\alpha=(p_y)_*[g_Ep_z^*\alpha],\qquad
 P_E\alpha=(p_y)_*[d'g_E\wedge p_z^*\alpha],\qquad
 Q_E\alpha=(p_y)_*[d''g_E\wedge p_z^*\alpha]
\tag{PB15.5}\label{r15b:rank-two-outputs-formula}
\]
have global ordinary PWS representatives of degrees
$(0,0),(1,0),(0,1)$. They have the prescribed pairing against
every ordinary smooth compact complementary test on $E$.
The proper current pushforward along $\pi$ of any PWS form
also has its ordinary PWS representative in the shifted degree.
\end{proposition}

\begin{proof}
The global singular expressions have the same diagonal-excision
meaning as in Round 14. The two pullbacks of $\alpha$ agree
on the diagonal, hence on an ordinary neighborhood of it by
the frozen immersion-germ theorem. Against the complementary
tests for the three displayed degrees, replacing $\alpha_z$
there by $\alpha_y$ makes at least one target degree exceed
$\dim E$. Thus a sufficiently small diagonal cutoff changes
none of the currents. Properness supplies such a cutoff on
each paired compact part, and globally here because $E$ is
projective.

On a trivializing open $U$ use
\eqref{r15b:rank-two-local-kernel}. The fixed-kernel part has
its prescribed outputs by the second part of
$(\mathrm{HP})_K$. The difference is a global PWS scalar on
$U\times\mathbf P^1_z\times\mathbf P^1_y$. Its products
with $\alpha_z$, $d'$ of that scalar with $\alpha_z$, and
$d''$ of that scalar with $\alpha_z$ are PWS of degrees
$(1,1),(2,1),(1,2)$. The first part of $(\mathrm{HP})_K$,
applied with target $U\times\mathbf P^1_y$, integrates all
three to ordinary PWS tensors with all mixed-test pairings.
No multiplication of an unspecified distribution by a PWS
tensor is being postulated.

The local sums therefore represent the global currents in
\eqref{r15b:rank-two-outputs-formula}. Actual-current
injectivity makes them equal on ordinary overlaps, so they
glue. The assertion for $\pi$ follows in the same way from
its local trivializations and the prescribed-current identity
in the first part of $(\mathrm{HP})_K$.
\end{proof}

\begin{theorem}[The Aeppli formula for a rank-two projective bundle]
\label{r15b:rank-two-aeppli}
Assume $(\mathrm{HP})_K$ and retain the frozen first-strip PWS
replacement, proper scalar constancy, and finite trace results.
For every $S,N,E$ as above,
\[
 H_{\mathrm A,\mathrm{act}}^{1,1}(E)
       =\pi^*H_{\mathrm A,\mathrm{act}}^{1,1}(S)
                                      \oplus\mathbf R[h].
\tag{PB15.6}\label{r15b:rank-two-aeppli-formula}
\]
The class $[h]$ is the actual class of any model first Chern
form for $\mathcal O_E(1)$. No computation of
$H_{\mathrm A}^{1,1}(S)$ is assumed.
\end{theorem}

\begin{proof}
The frozen replacement theorem supplies a $D$-closed PWS
representative $\alpha$ of an arbitrary actual Aeppli class
on $E$. Define ordinary PWS outputs
\[
                     I\alpha=\pi_*\alpha,\qquad
                     J\alpha=\pi_*(h\wedge\alpha).
\tag{PB15.7}\label{r15b:rank-two-IJ}
\]
They have degrees $(0,0)$ and $(1,1)$ on $S$, respectively.
The second input is PWS because $h$ is smooth. Proper-current
differentiation and actual Stokes give
\[
                             D(I\alpha)=D(J\alpha)=0.
\]
These are actual identities by current injectivity, using
$D\alpha=0$ and separate closedness of $h$. The frozen
projective scalar constancy theorem makes $I\alpha=c$ a
real constant.

Here is a direct all-smooth-test derivation of the needed
homotopy, which also fixes its signs. Put $m=\dim E$ and let
$\eta\in A_c^{m-1,m-1}(E)$ be arbitrary. Apply the scalar
Green identity \eqref{r15b:rank-two-Green} to the global actual
test $\Phi=\alpha_z\wedge p_y^*\eta$ on $T$, of degree
$(m,m)$. It gives
\[
 \int_Tg_E D\Phi
   =\int_T(h_z+h_y+\pi_T^*\kappa)\wedge\Phi
                                     -\int_E\alpha\wedge\eta.
\tag{PB15.8}\label{r15b:rank-two-scalar-pairing}
\]
On a neighborhood of the diagonal, the forms
$\alpha_z\wedge p_y^*d'\eta$ and
$\alpha_z\wedge p_y^*d''\eta$ vanish, since there one can
replace $\alpha_z$ by $\alpha_y$ and one of their base degrees
exceeds $m$. Hence their products with $g_E$ extend by zero
across the diagonal as actual compactly supported forms.
Applying actual Stokes to those two products, and using
$D\alpha=0$, gives
\[
 \begin{split}
 \big\langle DG_E\alpha+d''P_E\alpha-d'Q_E\alpha,\eta\big\rangle
 &=\int_T\bigl(g_E\alpha_z\wedge D\eta_y
      +d'g_E\wedge\alpha_z\wedge d''\eta_y
      -d''g_E\wedge\alpha_z\wedge d'\eta_y\bigr)\\
 &=-\int_Tg_E D(\alpha_z\wedge\eta_y).
 \end{split}
\tag{PB15.9}\label{r15b:rank-two-boundary-sign}
\]
For example
$D(\alpha_z\eta_y)=-d''\alpha_z\,d'\eta_y
 +d'\alpha_z\,d''\eta_y+\alpha_zD\eta_y$;
the two Stokes identities cancel these first-derivative terms
with precisely the signs in \eqref{r15b:rank-two-boundary-sign}.
Thus no differentiated cutoff term has been discarded.

The coordinate base-change clause of $(\mathrm{HP})_K$,
checked on local trivializations and then glued by current
injectivity, identifies
\[
 (p_y)_*\alpha_z=\pi^*I\alpha,
 \qquad (p_y)_*(h_z\alpha_z)=\pi^*J\alpha.
\]
The ordinary projection formula with the smooth target forms
$h$ and $\pi^*\kappa$ handles the other two curvature terms.
Equations \eqref{r15b:rank-two-scalar-pairing} and
\eqref{r15b:rank-two-boundary-sign} consequently prove, for
every prescribed smooth test, the identity
\[
 \alpha=c h+\pi^*(J\alpha+c\kappa)
                  +DG_E\alpha+d''P_E\alpha-d'Q_E\alpha.
\tag{PB15.10}\label{r15b:rank-two-homotopy}
\]
All terms are ordinary actual forms. Current injectivity makes
this an actual identity. The last three terms are Aeppli
boundaries, while $J\alpha+c\kappa$ is $D$-closed on $S$.
This proves spanning in \eqref{r15b:rank-two-aeppli-formula}.

For directness restrict a relation
$\pi^*[\beta]+a[h]=0$ to a fiber over a closed point $s$ of
$S$. The first term restricts to zero. The native mass of
$h$ on that projective-line fiber is $[\kappa(s):K]$:
compare its smooth metric with the standard model metric
using the global PWS ratio and actual Stokes on the fiber.
Therefore $a=0$. Finally $\pi^*$ is injective on actual
$H_{\mathrm A}^{1,1}$ by the frozen R14 constant-fiber-dimension
theorem, \texttt{r14j:equidimensional}, applied to the smooth
projective surjection $\pi$. This gives $[\beta]=0$.
The same PWS metric-ratio argument identifies $[h]$ with its
model first Chern class. Every step retains native degrees.
\end{proof}

\begin{corollary}[Codimension-two smooth centers in any ambient dimension]
\label{r15b:codim-two-center}
Assume $(\mathrm{HP})_K$. Let $X/K$ be smooth integral
projective and let $C\subset X$ be a smooth closed center of
pure codimension two. Then the blowup $q:Y\to X$ induces
an injection on actual $H_{\mathrm{BC}}^{2,2}$.
\end{corollary}

\begin{proof}
Treat the disjoint integral components of $C$ separately.
The exceptional divisor is $E=\mathbf P_C(N_{C/X})$ with
rank-two normal bundle. Theorem~\ref{r15b:rank-two-aeppli}
gives its full Aeppli group as the sum of a center pullback
and the relative hyperplane class. The latter is the negative
of $i^*c_1(\mathcal O_Y(E))$. Both summands therefore vanish
in the normalized exceptional quotient of
Proposition~\ref{r15:degree-two-all-centers}. That proposition,
which has no center-dimension restriction, identifies this
quotient with the kernel of $q^*$ in curvature degree $(2,2)$.
The quotient is zero, as required.
\end{proof}

\paragraph{Exact limits.}
This argument treats rank two. It does not automatically give
an Aeppli formula for $\mathbf P_S(N)$ when
$\operatorname{rank}N>2$: its relative diagonal is no longer
a divisor and the scalar kernel used here is unavailable.
It gives no general proper actual pushforward and no all-cycle
Green theorem in higher ambient dimension. It does remove the
surface-center analytic quotient in a fourfold, conditional
on $(\mathrm{HP})_K$, without computing the surface's own
Aeppli group. The hypothesis $(\mathrm{HP})_K$ is supplied
by Theorem~\ref{r15hp:projection} and its Green-output
consequence. The separate incidence theorem
\ref{r15i:aeppli} establishes the all-rank $(1,1)$
formula; it is not inferred merely from the rank-two
diagonal calculation above.

\paragraph{Primary metric references checked.}
CLD, \emph{Formes differentielles reelles et courants sur les
espaces de Berkovich}, 2025 version,
\url{https://arxiv.org/pdf/1204.6277}, Propositions 15.2.8,
17.1.7 and 17.1.8; Remark 17.1.9 distinguishes higher-rank
vector norms. The scalar actual divisor Green identity and
actual-current injectivity are the previously fixed Mihatsch
inputs. The first-strip PWS replacement is
\texttt{round14/projective\_cohomology/pws\_replacement\_degree\_11.tex},
and scalar constancy is
\texttt{round14/foundations\_audit/proper\_scalar\_constancy.tex}.
